% !TEX root = ../article.tex
% =============================================================================
% 3.6 Evaluación de la solución
% =============================================================================
% Guía: métricas y criterios adecuados para la tarea, con su justificación;
% comparación entre la solución base y la propuesta en las mismas
% condiciones; resultados con tablas, gráficos y ejemplos de aciertos y
% errores; tiempo de respuesta, consumo de recursos o costo cuando sean
% relevantes. Debe incluir casos representativos del uso previsto.
%
% Segundo avance (asesoría 2, 7-11 dic 2026): diseño de la solución,
% métricas y procedimiento de evaluación; solución base y primeros
% resultados.
% =============================================================================

\section{Evaluación de la solución}\label{sec:evaluation}

\subsection{Métricas y criterios}\label{sec:metrics}
\pendiente{métricas de segmentación elegidas y por qué son adecuadas}

\subsection{Comparación con la solución base}\label{sec:comparison}
\pendiente{tabla con la solución base y la propuesta, evaluadas sobre el mismo
conjunto de prueba}

% Plantilla de tabla (descomentar cuando haya resultados):
% \begin{table}[htbp]
%     \centering
%     \caption{Comparación entre la solución base y la propuesta.}\label{tab:comparison}
%     \begin{tabular}{lcc}
%         \toprule
%         Modelo & Métrica 1 & Métrica 2 \\
%         \midrule
%         Solución base & -- & -- \\
%         Propuesta     & -- & -- \\
%         \bottomrule
%     \end{tabular}
%     \articleSource{Elaboración propia.}
% \end{table}

\subsection{Ejemplos de aciertos y errores}\label{sec:examples}
\pendiente{casos representativos del uso previsto: cortes con la máscara
predicha frente a la de referencia, incluidos casos de fallo}

\subsection{Evaluación de la incertidumbre}\label{sec:uncertainty-eval}
La incertidumbre (Sección~\ref{sec:explainability}) se evalúa con el mismo
conjunto de prueba interno y con la validación externa, en tres pruebas:

\begin{enumerate}
    \item \textbf{Relación con el error por caso:} correlación de Spearman
          entre la puntuación de confianza de cada estudio y su Dice. Una
          correlación clara indica que la puntuación sirve para marcar los
          casos que requieren revisión.
    \item \textbf{Relación con el error por vóxel:} capacidad del mapa de
          incertidumbre para distinguir los vóxeles mal segmentados de los
          bien segmentados, medida con el área bajo la curva ROC.
    \item \textbf{Calibración:} diagrama de confiabilidad que compara las
          probabilidades del modelo con la proporción real de aciertos, porque
          las U-Net tienden a ser demasiado confiadas \parencite{mehrtash2020}.
\end{enumerate}

\pendiente{resultados de las tres pruebas, con la solución base y la
propuesta}

\subsection{Tiempo de respuesta y recursos}\label{sec:resources}
\pendiente{tiempo de inferencia por estudio, memoria y costo de uso}
